\section{A}

% achten (to pay attention to / to respect)
\begin{verbentry}{
  style = {default},
  toc = {achten (to pay attention to / to respect)},
  name = {achten},
  english = {to pay attention to / to respect},
  type = {regular},
  auxiliary = {haben}
}
 
    
     \vspace{5pt}

    \hrule
    
    \vspace{5pt}

  \begin{tabular}{rl}
     \multicolumn{2}{l}{\textbf{PRÄSENS} }\\
    ich & \textbf{achte}\\
    du & \textbf{achtest}\\
    er/sie/es & \textbf{achtet}\\
    wir & \textbf{achten}\\
    ihr & \textbf{achtet}\\
    sie/Sie & \textbf{achten}\\
  \end{tabular}
  \hfill
     \begin{tabular}{rl}
    \multicolumn{2}{l}{\textbf{PRÄTERITUM} }\\
    ich & \textbf{achtete}\\
    du & \textbf{achtetest}\\
    er/sie/es & \textbf{achtete}\\
    wir & \textbf{achteten}\\
    ihr & \textbf{achtetet}\\
    sie/Sie & \textbf{achteten}\\
  \end{tabular}
  \hfill
  \begin{tabular}{rl}
     \multicolumn{2}{l}{\textbf{IMPERATIV} }\\
   & \textbf{acht(e) (du)!}\\
   & \textbf{achten wir!}\\
   & \textbf{achtet (ihr)!}\\
   & \textbf{achten Sie!}\\
   & \vspace{10pt}
  \end{tabular}

    \vspace{10pt}

 \begin{tabular}{rl}
     \multicolumn{2}{l}{\textbf{PERFEKT}}\\
   & haben + \textbf{geachtet}\\
  \end{tabular}
\hfill
   \begin{tabular}{rl}
   
    \multicolumn{2}{l}{\textbf{FUTURE I} }\\
   & werden + \textbf{achten}\\
 
  \end{tabular}
\hfill
   \begin{tabular}{rl}
      \multicolumn{2}{l}{\textbf{PLUSQUAMPERFEKT}}\\
   & hatten + \textbf{geachtet}\\
  \end{tabular}

   \vspace{10pt}

   \begin{tabular}{lll}
      \multicolumn{3}{l}{\textbf{MIT PRÄPOSITIONEN}}\\
& \textbf{achten auf} + Akk & (to pay attention to)\\
& \textbf{achten} + Akk & (to respect / value)\\
\end{tabular}
  \hfill
   \begin{tabular}{lll}
\multicolumn{3}{l}{\textbf{TRENNBARE VERBEN}}\\
& \textbf{---} & \\
\end{tabular}


  \vspace{-5pt}
\end{verbentry}

% ahmen (to imitate)
\begin{verbentry}{
  style = {default},
  toc = {ahmen (to imitate)},
  name = {ahmen},
  english = {to imitate},
  type = {regular},
  auxiliary = {haben}
}


\vspace{5pt}
\hrule
\vspace{5pt}

\begin{tabular}{rl}
\multicolumn{2}{l}{\textbf{PRÄSENS}}\\
ich & \textbf{ahme}\\
du & \textbf{ahmst}\\
er/sie/es & \textbf{ahmt}\\
wir & \textbf{ahmen}\\
ihr & \textbf{ahmt}\\
sie/Sie & \textbf{ahmen}\\
\end{tabular}
\hfill
\begin{tabular}{rl}
\multicolumn{2}{l}{\textbf{PRÄTERITUM}}\\
ich & \textbf{ahmte}\\
du & \textbf{ahmtest}\\
er/sie/es & \textbf{ahmte}\\
wir & \textbf{ahmten}\\
ihr & \textbf{ahmtet}\\
sie/Sie & \textbf{ahmten}\\
\end{tabular}
\hfill
\begin{tabular}{rl}
\multicolumn{2}{l}{\textbf{IMPERATIV}}\\
& \textbf{ahme (du)!}\\
& \textbf{ahmen wir!}\\
& \textbf{ahmt (ihr)!}\\
& \textbf{ahmen Sie!}\\
\end{tabular}

\vspace{10pt}

\begin{tabular}{rl}
\multicolumn{2}{l}{\textbf{PERFEKT}}\\
& haben + \textbf{geahmt}\\
\end{tabular}
\hfill
\begin{tabular}{rl}
\multicolumn{2}{l}{\textbf{FUTURE I}}\\
& werden + \textbf{ahmen}\\
\end{tabular}
\hfill
\begin{tabular}{rl}
\multicolumn{2}{l}{\textbf{PLUSQUAMPERFEKT}}\\
& hatten + \textbf{geahmt}\\
\end{tabular}

\vspace{10pt}

\begin{tabular}{lll}
\multicolumn{3}{l}{\textbf{MIT PRÄPOSITIONEN}}\\
& \textbf{---} & \\
\end{tabular}
\hfill
\begin{tabular}{lll}
\multicolumn{3}{l}{\textbf{TRENNBARE VERBEN}}\\
& \textbf{nach$|$ahmen} & (to imitate / emulate)\\
\end{tabular}
\vspace{-5pt}
\end{verbentry}

% akzeptieren (to accept)
\begin{verbentry}{
  style = {default},
  toc = {akzeptieren (to accept)},
  name = {akzeptieren},
  english = {to accept},
  type = {regular},
  auxiliary = {haben}
}
 

\vspace{5pt}
\hrule
\vspace{5pt}

\begin{tabular}{rl}
\multicolumn{2}{l}{\textbf{PRÄSENS}}\\
ich & \textbf{akzeptiere}\\
du & \textbf{akzeptierst}\\
er/sie/es & \textbf{akzeptiert}\\
wir & \textbf{akzeptieren}\\
ihr & \textbf{akzeptiert}\\
sie/Sie & \textbf{akzeptieren}\\
\end{tabular}
\hfill
\begin{tabular}{rl}
\multicolumn{2}{l}{\textbf{PRÄTERITUM}}\\
ich & \textbf{akzeptierte}\\
du & \textbf{akzeptiertest}\\
er/sie/es & \textbf{akzeptierte}\\
wir & \textbf{akzeptierten}\\
ihr & \textbf{akzeptiertet}\\
sie/Sie & \textbf{akzeptierten}\\
\end{tabular}
\hfill
\begin{tabular}{rl}
\multicolumn{2}{l}{\textbf{IMPERATIV}}\\
& \textbf{akzeptiere (du)!}\\
& \textbf{akzeptieren wir!}\\
& \textbf{akzeptiert (ihr)!}\\
& \textbf{akzeptieren Sie!}\\
\end{tabular}

\vspace{10pt}

\begin{tabular}{rl}
\multicolumn{2}{l}{\textbf{PERFEKT}}\\
& haben + \textbf{akzeptiert}\\
\end{tabular}
\hfill
\begin{tabular}{rl}
\multicolumn{2}{l}{\textbf{FUTURE I}}\\
& werden + \textbf{akzeptieren}\\
\end{tabular}
\hfill
\begin{tabular}{rl}
\multicolumn{2}{l}{\textbf{PLUSQUAMPERFEKT}}\\
& hatten + \textbf{akzeptiert}\\
\end{tabular}

\vspace{10pt}

\begin{tabular}{lll}
\multicolumn{3}{l}{\textbf{MIT PRÄPOSITIONEN}}\\
& \textbf{akzeptieren} + Akk & (to accept something)\\
\end{tabular}
\hfill
\begin{tabular}{lll}
\multicolumn{3}{l}{\textbf{TRENNBARE VERBEN}}\\
& \textbf{---} & \\
\end{tabular}

\vspace{-5pt}
\end{verbentry}

% ändern (to change)
\begin{verbentry}{
  style = {acc},
  toc = {ändern (to change)},
  name = {ändern},
  english = {to change},
  type = {regular, + Akkusativ},
  auxiliary = {haben}
}
 
    
     \vspace{5pt}

    \hrule
    
    \vspace{5pt}

  \begin{tabular}{rl}
     \multicolumn{2}{l}{\textbf{PRÄSENS} }\\
    ich & \textbf{ändere}\\
    du & \textbf{änderst}\\
    er/sie/es & \textbf{ändert}\\
    wir & \textbf{ändern}\\
    ihr & \textbf{ändert}\\
    sie/Sie & \textbf{ändern}\\
  \end{tabular}
  \hfill
     \begin{tabular}{rl}
    \multicolumn{2}{l}{\textbf{PRÄTERITUM} }\\
    ich & \textbf{änderte}\\
    du & \textbf{ändertest}\\
    er/sie/es & \textbf{änderte}\\
    wir & \textbf{änderten}\\
    ihr & \textbf{ändertet}\\
    sie/Sie & \textbf{änderten}\\
  \end{tabular}
  \hfill
  \begin{tabular}{rl}
     \multicolumn{2}{l}{\textbf{IMPERATIV} }\\
   & \textbf{änder(e) (du)!}\\
   & \textbf{ändern wir!}\\
   & \textbf{ändert (ihr)!}\\
   & \textbf{ändern Sie!}\\
   & \vspace{10pt}
  \end{tabular}

    \vspace{10pt}

 \begin{tabular}{rl}
     \multicolumn{2}{l}{\textbf{PERFEKT}}\\
  &  haben + \textbf{geändert}\\
  \end{tabular}
\hfill
   \begin{tabular}{rl}
   
    \multicolumn{2}{l}{\textbf{FUTURE I} }\\
  &  werden + \textbf{ändern}\\
 
  \end{tabular}
\hfill
   \begin{tabular}{rl}
      \multicolumn{2}{l}{\textbf{PLUSQUAMPERFEKT}}\\
  &  hatten + \textbf{geändert}\\
  \end{tabular}
  

  \vspace{10pt}

\begin{tabular}{lll}
\multicolumn{3}{l}{\textbf{MIT PRÄPOSITIONEN}}\\
& \textbf{---} & \\
\end{tabular}
\hfill
\begin{tabular}{lll}
\multicolumn{3}{l}{\textbf{TRENNBARE VERBEN}}\\
& \textbf{---} & \\
\end{tabular}

  \vspace{-5pt}
\end{verbentry}

% antworten (to answer)
\begin{verbentry}{
  style = {dat},
  toc = {antworten (to answer)},
  name = {antworten},
  english = {to answer},
  type = {regular, + Dativ},
  auxiliary = {haben}
}
 
    
     \vspace{5pt}

    \hrule
    
    \vspace{5pt}

  \begin{tabular}{rl}
     \multicolumn{2}{l}{\textbf{PRÄSENS} }\\
    ich & \textbf{antworte}\\
    du & \textbf{antwortest}\\
    er/sie/es & \textbf{antwortet}\\
    wir & \textbf{antworten}\\
    ihr & \textbf{antwortet}\\
    sie/Sie & \textbf{antworten}\\
  \end{tabular}
  \hfill
     \begin{tabular}{rl}
    \multicolumn{2}{l}{\textbf{PRÄTERITUM} }\\
    ich & \textbf{antwortete}\\
    du & \textbf{antwortetest}\\
    er/sie/es & \textbf{antwortete}\\
    wir & \textbf{antworteten}\\
    ihr & \textbf{antwortetet}\\
    sie/Sie & \textbf{antworteten}\\
  \end{tabular}
  \hfill
  \begin{tabular}{rl}
     \multicolumn{2}{l}{\textbf{IMPERATIV} }\\
   & \textbf{antwort(e) (du)!}\\
   & \textbf{antworten wir!}\\
   & \textbf{antwortet (ihr)!}\\
   & \textbf{antworten Sie!}\\
   & \vspace{10pt}
  \end{tabular}

    \vspace{10pt}

 \begin{tabular}{rl}
     \multicolumn{2}{l}{\textbf{PERFEKT}}\\
  &  haben + \textbf{geantwortet}\\
  \end{tabular}
\hfill
   \begin{tabular}{rl}
   
    \multicolumn{2}{l}{\textbf{FUTURE I} }\\
   & werden + \textbf{antworten}\\
 
  \end{tabular}
\hfill
   \begin{tabular}{rl}
      \multicolumn{2}{l}{\textbf{PLUSQUAMPERFEKT}}\\
   & hatten + \textbf{geantwortet}\\
  \end{tabular}


    \vspace{10pt}

\begin{tabular}{lll}
\multicolumn{3}{l}{\textbf{MIT PRÄPOSITIONEN}}\\
& \textbf{antworten auf} + Akk & (to answer something)\\
\end{tabular}
  \hfill
\begin{tabular}{lll}
\multicolumn{3}{l}{\textbf{TRENNBARE VERBEN}}\\
& \textbf{---} & \\
\end{tabular}

  \vspace{-5pt}
\end{verbentry}

% arbeiten (to work)
\begin{verbentry}{
  style = {default},
  toc = {arbeiten (to work)},
  name = {arbeiten},
  english = {to work},
  type = {regular},
  auxiliary = {haben}
}
 
    
     \vspace{5pt}

    \hrule
    
    \vspace{5pt}

  \begin{tabular}{rl}
     \multicolumn{2}{l}{\textbf{PRÄSENS} }\\
    ich & \textbf{arbeite}\\
    du & \textbf{arbeitest}\\
    er/sie/es & \textbf{arbeitet}\\
    wir & \textbf{arbeiten}\\
    ihr & \textbf{arbeitet}\\
    sie/Sie & \textbf{arbeiten}\\
  \end{tabular}
  \hfill
     \begin{tabular}{rl}
    \multicolumn{2}{l}{\textbf{PRÄTERITUM} }\\
    ich & \textbf{arbeitete}\\
    du & \textbf{arbeitetest}\\
    er/sie/es & \textbf{arbeitete}\\
    wir & \textbf{arbeiteten}\\
    ihr & \textbf{arbeitetet}\\
    sie/Sie & \textbf{arbeiteten}\\
  \end{tabular}
  \hfill
  \begin{tabular}{rl}
     \multicolumn{2}{l}{\textbf{IMPERATIV} }\\
   & \textbf{arbeit(e) (du)!}\\
   & \textbf{arbeiten wir!}\\
   & \textbf{arbeitet (ihr)!}\\
   & \textbf{arbeiten Sie!}\\
   & \vspace{10pt}
  \end{tabular}

    \vspace{10pt}

 \begin{tabular}{rl}
     \multicolumn{2}{l}{\textbf{PERFEKT}}\\
   & haben + \textbf{gearbeitet}\\
  \end{tabular}
\hfill
   \begin{tabular}{rl}
   
    \multicolumn{2}{l}{\textbf{FUTURE I} }\\
   & werden + \textbf{arbeiten}\\
 
  \end{tabular}
\hfill
   \begin{tabular}{rl}
      \multicolumn{2}{l}{\textbf{PLUSQUAMPERFEKT}}\\
   & hatten + \textbf{gearbeitet}\\
  \end{tabular}
  

    \vspace{10pt}
    
   \begin{tabular}{lll}
      \multicolumn{3}{l}{\textbf{MIT PRÄPOSITIONEN}}\\
& \textbf{arbeiten an} + Dat & (to work on)\\
& \textbf{arbeiten für} + Akk & (to work for)\\
   \end{tabular}
  \hfill
   \begin{tabular}{lll}
      \multicolumn{3}{l}{\textbf{TRENNBARE VERBEN}}\\
\vspace{20pt}
  \end{tabular}


  \vspace{-5pt}
\end{verbentry}

% ärgern (to annoy)
\begin{verbentry}{
  style = {default},
  toc = {ärgern (to annoy)},
  name = {ärgern},
  english = {to annoy},
  type = {regular},
  auxiliary = {haben}
}
 
    
     \vspace{5pt}

    \hrule
    
    \vspace{5pt}

  \begin{tabular}{rl}
     \multicolumn{2}{l}{\textbf{PRÄSENS} }\\
    ich & \textbf{ärgere}\\
    du & \textbf{ärgerst}\\
    er/sie/es & \textbf{ärgert}\\
    wir & \textbf{ärgern}\\
    ihr & \textbf{ärgert}\\
    sie/Sie & \textbf{ärgern}\\
  \end{tabular}
  \hfill
     \begin{tabular}{rl}
    \multicolumn{2}{l}{\textbf{PRÄTERITUM} }\\
    ich & \textbf{ärgerte}\\
    du & \textbf{ärgertest}\\
    er/sie/es & \textbf{ärgerte}\\
    wir & \textbf{ärgerten}\\
    ihr & \textbf{ärgertet}\\
    sie/Sie & \textbf{ärgerten}\\
  \end{tabular}
  \hfill
  \begin{tabular}{rl}
     \multicolumn{2}{l}{\textbf{IMPERATIV} }\\
   & \textbf{ärger(e) (du)!}\\
   & \textbf{ärgern wir!}\\
   & \textbf{ärgert (ihr)!}\\
   & \textbf{ärgern Sie!}\\
   & \vspace{10pt}
  \end{tabular}

    \vspace{10pt}

 \begin{tabular}{rl}
     \multicolumn{2}{l}{\textbf{PERFEKT}}\\
   & haben + \textbf{geärgert}\\
  \end{tabular}
\hfill
   \begin{tabular}{rl}
   
    \multicolumn{2}{l}{\textbf{FUTURE I} }\\
   & werden + \textbf{ärgern}\\
 
  \end{tabular}
\hfill
   \begin{tabular}{rl}
      \multicolumn{2}{l}{\textbf{PLUSQUAMPERFEKT}}\\
   & hatten + \textbf{geärgert}\\
  \end{tabular}

   \vspace{10pt}

   \begin{tabular}{lll}
      \multicolumn{3}{l}{\textbf{MIT PRÄPOSITIONEN}}\\
& \textbf{ärgern über} + Akk & (to be annoyed about)\\
   \end{tabular}
  \hfill
   \begin{tabular}{lll}
      \multicolumn{3}{l}{\textbf{TRENNBARE VERBEN}}\\
\vspace{10pt}
  \end{tabular}

\vspace{-5pt}
\end{verbentry}

% atmen (to breathe)
\begin{verbentry}{
  style = {default},
  toc = {atmen (to breathe)},
  name = {atmen},
  english = {to breathe},
  type = {regular},
  auxiliary = {haben}
}


\vspace{5pt}
\hrule
\vspace{5pt}

\begin{tabular}{rl}
\multicolumn{2}{l}{\textbf{PRÄSENS}}\\
ich & \textbf{atme}\\
du & \textbf{atmest}\\
er/sie/es & \textbf{atmet}\\
wir & \textbf{atmen}\\
ihr & \textbf{atmet}\\
sie/Sie & \textbf{atmen}\\
\end{tabular}
\hfill
\begin{tabular}{rl}
\multicolumn{2}{l}{\textbf{PRÄTERITUM}}\\
ich & \textbf{atmete}\\
du & \textbf{atmetest}\\
er/sie/es & \textbf{atmete}\\
wir & \textbf{atmeten}\\
ihr & \textbf{atmetet}\\
sie/Sie & \textbf{atmeten}\\
\end{tabular}
\hfill
\begin{tabular}{rl}
\multicolumn{2}{l}{\textbf{IMPERATIV}}\\
& \textbf{atme (du)!}\\
& \textbf{atmen wir!}\\
& \textbf{atmet (ihr)!}\\
& \textbf{atmen Sie!}\\
\end{tabular}

\vspace{10pt}

\begin{tabular}{rl}
\multicolumn{2}{l}{\textbf{PERFEKT}}\\
& haben + \textbf{geatmet}\\
\end{tabular}
\hfill
\begin{tabular}{rl}
\multicolumn{2}{l}{\textbf{FUTURE I}}\\
& werden + \textbf{atmen}\\
\end{tabular}
\hfill
\begin{tabular}{rl}
\multicolumn{2}{l}{\textbf{PLUSQUAMPERFEKT}}\\
& hatten + \textbf{geatmet}\\
\end{tabular}

\vspace{10pt}

\begin{tabular}{lll}
\multicolumn{3}{l}{\textbf{MIT PRÄPOSITIONEN}}\\
& \textbf{atmen in} + Akk & (to breathe into)\\
& \textbf{atmen durch} + Akk & (to breathe through)\\
& \textbf{atmen mit} + Dat & (to breathe with)\\
\end{tabular}
\hfill
\begin{tabular}{lll}
\multicolumn{3}{l}{\textbf{TRENNBARE VERBEN}}\\
& \textbf{ein$|$atmen} & (to inhale)\\
& \textbf{aus$|$atmen} & (to exhale)\\
& \textbf{auf$|$atmen} & (to breathe again / breathe easier)\\
\end{tabular}
\vspace{-5pt}
\end{verbentry}
